\documentclass[UTF8,11pt,fontset=none]{ctexart}
\usepackage[a4paper,margin=23mm,headheight=15pt,footskip=11mm]{geometry}
\usepackage{fontspec}
\IfFontExistsTF{TeX Gyre Termes}{\setmainfont{TeX Gyre Termes}}{\setmainfont{lmroman10-regular.otf}[BoldFont=lmroman10-bold.otf,ItalicFont=lmroman10-italic.otf,BoldItalicFont=lmroman10-bolditalic.otf]}
\IfFontExistsTF{TeX Gyre Heros}{\setsansfont{TeX Gyre Heros}}{\setsansfont{lmsans10-regular.otf}[BoldFont=lmsans10-bold.otf,ItalicFont=lmsans10-oblique.otf]}
\IfFontExistsTF{DejaVu Sans Mono}{\setmonofont{DejaVu Sans Mono}[Scale=MatchLowercase]}{\setmonofont{lmmono10-regular.otf}}
\IfFontExistsTF{Noto Serif CJK SC}{\setCJKmainfont{Noto Serif CJK SC}[AutoFakeSlant=0.15]}{\setCJKmainfont{FandolSong-Regular.otf}[BoldFont=FandolSong-Bold.otf,ItalicFont=FandolKai-Regular.otf]}
\IfFontExistsTF{Noto Sans CJK SC}{\setCJKsansfont{Noto Sans CJK SC}}{\setCJKsansfont{FandolHei-Regular.otf}}
\IfFontExistsTF{Noto Sans Mono CJK SC}{\setCJKmonofont{Noto Sans Mono CJK SC}}{\setCJKmonofont{FandolFang-Regular.otf}}
\usepackage{amsmath,amssymb,amsthm,mathtools,mathrsfs,esint,bm}
\usepackage{booktabs,longtable,array,calc}
\usepackage{enumitem}
\setlist{nosep,leftmargin=1.7em,topsep=3pt}
\usepackage{fvextra}
\DefineVerbatimEnvironment{Highlighting}{Verbatim}{breaklines,breakanywhere,fontsize=\small}
\usepackage{xurl}
\usepackage[unicode,hidelinks,bookmarksnumbered]{hyperref}
\usepackage{bookmark}
\usepackage{fancyhdr}
\pagestyle{fancy}\fancyhf{}
\fancyhead[L]{\small 第一科学主线：近期研究全稿}
\fancyhead[R]{\small 2026-10-07}\fancyfoot[C]{\thepage}
\renewcommand{\headrulewidth}{0.3pt}
\providecommand{\tightlist}{\setlength{\itemsep}{1pt}\setlength{\parskip}{0pt}}
\providecommand{\passthrough}[1]{#1}
\newcommand{\Ric}{\operatorname{Ric}}
\newcommand{\Vol}{\operatorname{Vol}}
\newcommand{\Rm}{\operatorname{Rm}}
\newcommand{\supp}{\operatorname{supp}}
\newcommand{\Var}{\operatorname{Var}}
\newcommand{\tr}{\operatorname{tr}}
\newcommand{\avint}{\fint}
\newcommand{\Jcal}{\mathcal J}
\newcommand{\Ecal}{\mathcal E}
\newcommand{\dd}{\,\mathrm d}
\newcommand{\Rea}{\operatorname{Re}}
\setlength{\parindent}{2em}
\setlength{\parskip}{3pt plus 1pt}
\setlength{\emergencystretch}{3em}
\linespread{1.10}\allowdisplaybreaks[2]
\setlength{\abovedisplayskip}{8pt plus 2pt minus 2pt}
\setlength{\belowdisplayskip}{8pt plus 2pt minus 2pt}
\renewcommand{\arraystretch}{1.12}
\setcounter{tocdepth}{1}\setcounter{secnumdepth}{2}
\pdfstringdefDisableCommands{\def\fint{int}\def\mathcal#1{#1}\def\mathfrak#1{#1}\def\mathscr#1{#1}\def\operatorname#1{#1}\def\varepsilon{epsilon}\def\chi{chi}\def\beta{beta}\def\rho{rho}\def\tau{tau}\def\Phi{Phi}\def\sigma{sigma}\def\kappa{kappa}\def\nabla{D}\def\mathbb#1{#1}\def\({}\def\){}\def\textsuperscript#1{#1}\def\textsubscript#1{#1}}
\title{同丛有界截面与第二陈类的曲率集中}
\author{第一科学主线·近期研究分卷}
\date{2026年10月7日}
\begin{document}
\maketitle
\noindent\textbf{原稿编号：}R02。

本章使用外部网页补充的强流结论 R1--R2，构造同一负手性丛中的有界替代截面，并以第二陈类检测曲率集中。保留替代截面的完整构造、慢梯度估计、陈类局部化和密度阈值推论。\cite{R02,Strong}

\section{实际同解输入}\label{ch2:ux5b9eux9645ux540cux89e3ux8f93ux5165}

考虑一列闭、定向、连通四维流形及其度量。采用原稿已经证明的共形归约，令

\[
 s_{g_j}=-4,\qquad D_{A_j}\Phi_j=0,\qquad
 F_{A_j}^{+}=2i\sigma(\Phi_j),\qquad 0\le\rho_j=|\Phi_j|^2\le2.
\]

原实际供给与拓扑假设保持为

\[
 p_j=(a_{g_j}^{+})^2\ge A_j>0,\qquad a_j^2\ge T_{+,j}=2\chi_j+3\tau_j.
\]

下文 \(\tau_j\) 用于三次扭率截面时写作 \texttt{\textbackslash{}tau\_j\^{}\{\textbackslash{}mathrm\{sp\}\}}，而签名写 \texttt{\textbackslash{}operatorname\{sign\}(M\_j)}，避免混淆。连接符号 \(A_j\) 与拓扑预算 \(A_j\) 在源稿中重名；正文出现 \(F_A\) 时指连接，单独的 \(A_j\) 在拓扑比值中指预算。

设

\[
 \eta_j=\fint|\operatorname{Ric}_{g_j}+g_j|^2
       +D_{\beta_j}/V_j+\beta_j^{-2}\longrightarrow0,
 \qquad\beta_j\longrightarrow\infty.
\]

所有平均均相对于概率测度 \(d\mu_j=dV_j/V_j\)；子集积分也除以整个流形体积。

定义复秩二丛

\[
 E_j=S^-_j\otimes L_j,
\]

其 determinant 为 \(L_j^3\)，诱导连接记 \(\nabla_j\)，相应负手性 Dirac 算子记 \(\mathcal D_j=D^-_{3A_j}\)。实际三次扭率为

\[
 t_j:=\tau_j^{\mathrm{sp}}\in\Gamma(E_j),\quad
 m_j=|t_j|^2,\quad q_j=|\nabla_{A_j}\Phi_j|^2,
 \quad \rho_*=4/3,\quad m_*=128/81.
\]

继承最新 \path{01_O10_Strong_Current.md} 中的结论：

\[
 \begin{gathered}
 \|\mathcal D_jt_j\|_{2,\mathrm{av}}\to0,\quad \|m_j-m_*\|_{2,\mathrm{av}}\to0,\quad \|q_j-4/9\|_{2,\mathrm{av}}\to0,\\
 \rho_j\to\rho_*\quad\text{in normalized measure}.
 \end{gathered}
 \tag{2.1}\label{eq:r2:1}
\]

还继承同一实际解的曲率拟合：令

\[
 f_j=-iF_{A_j}^-,\quad Q_j^{\mathrm{sp}}=-i q_-(t_j),\quad
 S_j=m_*+2\rho_j^2q_j,\quad
 K_j=-\frac98\frac{\rho_j^3Q_j^{\mathrm{sp}}}{S_j}.
\]

则

\[
 |K_j|\le C,\qquad
 \|f_j-K_j\|_{2,\mathrm{av}}\to0,
 \qquad \|\operatorname{Ric}_{g_j}+g_j\|_{2,\mathrm{av}}\to0.
 \tag{2.2}\label{eq:r2:2}
\]

原稿 Clifford 规范是

\[
 c(q_-(t))t=\frac{|t|^2}{2}t.
 \tag{2.3}\label{eq:r2:3}
\]

这里没有增加 \(W^-\) 的界。后面需要有界总能量时会明确说明；若尾能量为无穷，相应必要下界自动成立。

\section{同丛替代截面}\label{ch2:ux540cux4e1bux66ffux4ee3ux622aux9762}

\subsection{命题 2.1：存在有界截面和有界 rough Laplacian 的同丛替代}\label{ch2:ux547dux9898-2.1ux5b58ux5728ux6709ux754cux622aux9762ux548cux6709ux754c-rough-laplacian-ux7684ux540cux4e1bux66ffux4ee3}

在 \eqref{eq:r2:1}--\eqref{eq:r2:3} 下，存在光滑 \(u_j\in\Gamma(E_j)\)，满足

\[
 \|u_j\|_\infty+\|\nabla_j^*\nabla_j u_j\|_\infty\le C,
 \tag{2.4}\label{eq:r2:4}
\]

以及

\[
 \|u_j-t_j\|_{2,\mathrm{av}}+
 \|\nabla_j(u_j-t_j)\|_{2,\mathrm{av}}\longrightarrow0.
 \tag{2.5}\label{eq:r2:5}
\]

特别地，\(|u_j|^2\to m_*\) 在归一化测度意义成立。\(u_j\) 不是另一个 SW 解，也不被宣称等于原来的三次扭率；它只作为同一丛的特征类测试截面使用。

\subsection{证明}\label{ch2:ux8bc1ux660e}

负手性 twisted Lichnerowicz 恒等式为

\[
 \mathcal D_j^*\mathcal D_j
 =\nabla_j^*\nabla_j-1+\frac32c(F_{A_j}^-).
\]

由 \eqref{eq:r2:3} 及 K 的定义，点态有

\[
 \frac32c(iK_j)t_j=-\gamma_jt_j,
 \qquad\gamma_j=\frac{27}{32}\frac{\rho_j^3m_j}{S_j}.
 \tag{2.6}\label{eq:r2:6}
\]

因为 \(m_j\le2\rho_j^2q_j\le S_j\)，所以 \(0\le\gamma_j\le27/4\)。由 \eqref{eq:r2:1}，

\[
 \gamma_j\longrightarrow1\text{ in measure}.
\]

置 \texttt{e\_j=(3/2)c(F\_\{A\_j\}\^{}-\ -\ iK\_j)}，则 \(\|e_j\|_{2,\rm av}\to0\)，且

\[
 (\nabla_j^*\nabla_j-2)t_j
 =\mathcal D_j^*\mathcal D_jt_j+(\gamma_j-1)t_j-e_jt_j.
 \tag{2.7}\label{eq:r2:7}
\]

取一个固定、光滑、纤维径向的截断 \texttt{T}，在 \texttt{\textbar{}t\textbar{}\textless{}=2} 上恒等，且 \texttt{\textbar{}T(t)\textbar{}\textless{}=4}。它由 Hermitian 范数定义，所以与丛的酉转移函数兼容。记 \texttt{z\_j=T(t\_j)}。由于 \(m_*<4\)、\texttt{m\_j\ -\textgreater{}\ m\_*} in L2av，

\[
 \|z_j-t_j\|_{2,\mathrm{av}}\longrightarrow0,
 \qquad \sup_j\|t_j\|_{4,\mathrm{av}}<\infty.
 \tag{2.8}\label{eq:r2:8}
\]

在每个闭流形上精确求解

\[
 \boxed{(4+\nabla_j^*\nabla_j)u_j=6z_j.}
 \tag{2.9}\label{eq:r2:9}
\]

正零阶项保证算子可逆；光滑性来自每个固定 j 上的椭圆理论。这里不需要几何一致的 Schauder 常数。点态最大值论证给

\[
 |u_j|\le6,\qquad
 |\nabla_j^*\nabla_j u_j|=|6z_j-4u_j|\le48.
\]

更具体地，在 \texttt{\textbar{}u\textbar{}\^{}2} 的最大点，采用非负 rough Laplacian 的规范，

\[
 0\le\Delta|u|^2
 =2\operatorname{Re}\langle6z-4u,u\rangle-2|\nabla u|^2,
\]

从而 \(4|u|\le6|z|\)（最大值非零时）。

令 \texttt{w\_j=u\_j-t\_j}、\texttt{d\_j=\textbackslash{}mathcal\ D\_jt\_j}。由 \eqref{eq:r2:7}、\eqref{eq:r2:9}，

\[
 (4+\nabla_j^*\nabla_j)w_j
 =6(z_j-t_j)-\mathcal D_j^*d_j-(\gamma_j-1)t_j+e_jt_j.
 \tag{2.10}\label{eq:r2:10}
\]

用 \(w_j\) 测试，在归一化体积下积分：

\[
 4\|w_j\|_2^2+\|\nabla_jw_j\|_2^2
 =6\langle z_j-t_j,w_j\rangle
 -\langle d_j,\mathcal D_jw_j\rangle
 -\langle(\gamma_j-1)t_j,w_j\rangle
 +\langle e_jt_j,w_j\rangle.
\]

此处及下一段的范数均为归一化范数。第一项用 \eqref{eq:r2:8} 和 Young；第二项用 \(|Dw|\le2|\nabla w|\) 和 \(\|d_j\|_2\to0\) 吸收。第三项趋零，因为

\[
 \|(\gamma_j-1)t_j\|_{4/3}
 \le\|\gamma_j-1\|_2\|t_j\|_4\to0.
\]

第四项满足

\[
 |\langle e_jt_j,w_j\rangle|
 \le\|e_j\|_2\|t_j\|_4\|w_j\|_4\to0.
\]

这里 \(w_j\) 的 \(L^4\) 界直接来自 \(u_j\) 的点态界与 \(t_j\) 的 \(L^4\) 界。吸收得到 \eqref{eq:r2:5}。证毕。

\section{有界 rough-Laplacian 版本的慢梯度截断}\label{ch2:ux6709ux754c-rough-laplacian-ux7248ux672cux7684ux6162ux68afux5ea6ux622aux65ad}

这是对所给 R1 证明结构的扩展。原 R1 用 \texttt{Delta\ u=f(\textbar{}u\textbar{}\^{}2)u}；此处只要求 \texttt{\textbar{}u\textbar{}+\textbar{}Delta\ u\textbar{}} 有界，不要求微分右端。

\subsection{引理 3.1：梯度高幅度能量由连接曲率的高幅度能量控制}\label{ch2:ux5f15ux7406-3.1ux68afux5ea6ux9ad8ux5e45ux5ea6ux80fdux91cfux7531ux8fdeux63a5ux66f2ux7387ux7684ux9ad8ux5e45ux5ea6ux80fdux91cfux63a7ux5236}

设 \texttt{E\_j} 是闭四流形上的固定秩 Hermitian 丛，连接曲率为 \texttt{F\_j}，且

\[
 |u_j|+|\nabla_j^*\nabla_j u_j|\le C_0,\qquad
 \operatorname{Ric}_{g_j}=R_{0,j}+r_j,\quad |R_{0,j}|\le C_0,\quad
 \|r_j\|_{2,\mathrm{av}}\to0.
 \tag{2.11}\label{eq:r2:11}
\]

先假设 \(\sup_j\fint|F_j|^2<\infty\)。定义

\[
 q=|\nabla u|^2,\quad H=\nabla^2u,\quad
 w_L=\left(\frac q{L^2+q}\right)^\epsilon.
\]

存在只依赖维数、秩与 \texttt{C\_0} 的固定 \(\epsilon>0\) 和常数 C，使对任意曲率截断值 \texttt{T\textgreater{}=1}，

\[
 \boxed{
 \fint w_L(|H|^2+q^2)
 \le C_T\,L^{-2\epsilon}
 +C\fint_{\{|F_j|>T\}}|F_j|^2
 +C\fint|r_j|^2.}
 \tag{2.12}\label{eq:r2:12}
\]

\(C_T\) 可依赖总连接曲率平均的上界；\textbf{乘在尾部能量前的 C 不依赖 T 或这个总能量上界。} 这一区别用于后面的普适浓缩门槛。

\subsection{证明}\label{ch2:ux8bc1ux660e-1}

令 \texttt{h=\textbackslash{}nabla\^{}*\textbackslash{}nabla\ u}，使用光滑权重

\[
 w_{L,\delta}=\left(\frac{q+\delta}{L^2+q+\delta}\right)^\epsilon,
 \qquad |dw|\sqrt q\le2\epsilon w|H|.
\]

设 \texttt{P\_i=\textbackslash{}nabla\_i\ u}，\texttt{H\_\{ij\}=\textbackslash{}nabla\_iP\_j}。保留曲率作用的 Hessian 换指标恒等式：

\[
 \int w|H|^2=
 \int w\operatorname{Re}\langle H_{ij},H_{ji}\rangle
 +\frac12\int w\sum_{i,j}|F_{ij}u|^2.
 \tag{2.13}\label{eq:r2:13}
\]

对交换项先积分掉一个导数，再交换 \texttt{\textbackslash{}nabla\_j\textbackslash{}nabla\_iP\_j}。协向量槽的曲率只给 Ricci 收缩，丛槽给 \texttt{F*P}。出现的 \texttt{\textbackslash{}nabla\ h} 项\textbf{再积分一次}，改写为 \texttt{\textbackslash{}int\ w\textbar{}h\textbar{}\^{}2} 以及 \texttt{dw*h*P}，因此不需要 \texttt{\textbackslash{}nabla\ h} 或曲率导数。取 epsilon 充分小，权重项吸收到左边，得到

\[
 I:=\int w|H|^2
 \le C\int w|h|^2
 +C\int w(|\operatorname{Ric}|+|F|)q
 +C\int w|F|^2|u|^2.
 \tag{2.14}\label{eq:r2:14}
\]

有界截面的加权插值来自 \texttt{div(w\ q\ Re\textless{}\ u,nabla\ u\ \textgreater{})} 的积分：

\[
 J:=\int wq^2
 \le C\int wq+C\sqrt{IJ}
 \le C_1\left(I+\int wq\right).
 \tag{2.15}\label{eq:r2:15}
\]

将 \texttt{F} 分成幅度不超过 T 的点态部分和尾部，将 Ricci 分成 \texttt{R\_0+r}；这些分解不被微分，不需要它们光滑。对尾部与 q 的乘积用

\[
 \int w(|F|1_{|F|>T}+|r|)q
 \le aJ+C_a\int(|F|^2 1_{|F|>T}+|r|^2).
\]

先固定 \texttt{C\_1}，再选 a 并吸收，得

\[
 I+J\le C_T\int w(1+q)
 +C\int(|F|^2 1_{|F|>T}+|r|^2).
 \tag{2.16}\label{eq:r2:16}
\]

这里尾部系数 C 的选择不依赖 T。用 \texttt{w=1} 及

\[
 \int q=\int\operatorname{Re}\langle h,u\rangle\le C V
\]

先取得

\[
 \fint(|H|^2+q^2)\le C\left(1+\fint|F|^2+\fint|r|^2\right).
\]

再用 \texttt{w\_\{L,delta\}\ \textless{}=\ L\^{}\{-2epsilon\}(q+delta)\^{}epsilon}、\texttt{0\textless{}epsilon\textless{}=1} 控制 \eqref{eq:r2:16} 中的加权低阶项。令 delta 降至零，得到 \eqref{eq:r2:12}。证毕。

\subsection{推论 3.2：任意小体积集合上的集中控制}\label{ch2:ux63a8ux8bba-3.2ux4efbux610fux5c0fux4f53ux79efux96c6ux5408ux4e0aux7684ux96c6ux4e2dux63a7ux5236}

对满足 \(\operatorname{Vol}(Z_j)/V_j\to0\) 的任意可测集合，定义

\[
 \mathfrak C_F:=\lim_{T\to\infty}\limsup_j
 \fint_{\{|F_j|>T\}}|F_j|^2.
\]

则

\[
 \boxed{
 \limsup_j\fint_{Z_j}|\nabla u_j|^4\le C\mathfrak C_F.}
 \tag{2.17}\label{eq:r2:17}
\]

证明中先固定 T,L。低梯度集合的贡献不超过 \(L^4\operatorname{Vol}(Z_j)/V_j\)；高梯度集合上 \(w_L\ge2^{-\epsilon}\)。先令 j 趋无穷，再令 L 趋无穷，最后令 T 趋无穷。此极限顺序消除 \(C_T\)，而保留 T 无关的尾部系数。特别地，连接曲率平方的一致可积性推出这些小集合上的梯度四次能量趋零。

\section{同丛特征类：把未估计的实际扭率通量换成可控制的测试}\label{ch2:ux540cux4e1bux7279ux5f81ux7c7bux628aux672aux4f30ux8ba1ux7684ux5b9eux9645ux626dux7387ux901aux91cfux6362ux6210ux53efux63a7ux5236ux7684ux6d4bux8bd5}

\subsection{命题 4.1：几乎定长的有界替代截面只能在曲率尾部承载第二 Chern 类}\label{ch2:ux547dux9898-4.1ux51e0ux4e4eux5b9aux957fux7684ux6709ux754cux66ffux4ee3ux622aux9762ux53eaux80fdux5728ux66f2ux7387ux5c3eux90e8ux627fux8f7dux7b2cux4e8c-chern-ux7c7b}

在引理 3.1 的条件下，设丛为复秩二，且

\[
 |u_j|^2\longrightarrow m_*>0\text{ in normalized measure}.
\]

则

\[
 \boxed{
 \limsup_j\frac{|c_2(E_j)[M_j]|}{V_j}
 \le C\mathfrak C_F.}
 \tag{2.18}\label{eq:r2:18}
\]

\subsection{证明}\label{ch2:ux8bc1ux660e-2}

固定 \texttt{0\textless{}a\textless{}b\textless{}m\_*}，取光滑标量 cutoff \texttt{chi}，在 \texttt{{[}0,a{]}} 上为1，在 \texttt{{[}b,infinity)} 上为0。令

\[
 Z_j=\{|u_j|^2\le b\},\qquad \operatorname{Vol}(Z_j)/V_j\to0.
\]

在 \(u_j\ne0\) 区域设 \texttt{n\_j=u\_j/\textbar{}u\_j\textbar{}}。用酉连接修正使 \texttt{n\_j} 平行，得到分裂为平凡复线与复线的连接 \(\nabla^0\)，其 \(c_2(F^0)=0\)。连接差 C 满足

\[
 |C|\le C|\nabla n_j|,\quad
 |F^t|\le C(|F_j|+|\nabla n_j|^2),
\]

于是通常的 Chern--Weil transgression 形式满足

\[
 dT_{n_j}=c_2(F_j),\qquad
 |T_{n_j}|\le C\left(|F_j||\nabla n_j|+|\nabla n_j|^3\right).
 \tag{2.19}\label{eq:r2:19}
\]

这是原整合稿已经使用的单位截面 transgression 构造，现用于 E 而非仅 S+。\((1-\chi(|u_j|^2))T_{n_j}\) 在零点附近恒零，可光滑延拓。Stokes 给出

\[
 c_2(E_j)[M_j]
 =\int\chi(|u_j|^2)c_2(F_j)
 +\int d\chi(|u_j|^2)\wedge T_{n_j}.
 \tag{2.20}\label{eq:r2:20}
\]

在过渡区 \(a\le|u_j|^2\le b\)，所有分母一致离零；\texttt{\textbar{}u\_j\textbar{}} 有界。因此

\[
 \frac{|c_2(E_j)[M_j]|}{V_j}
 \le C\fint_{Z_j}\left(|F_j|^2+|\nabla u_j|^4\right).
 \tag{2.21}\label{eq:r2:21}
\]

曲率项可按任意固定 T 分成 \(T^2\operatorname{Vol}(Z_j)/V_j\) 与尾部。梯度项由 \eqref{eq:r2:17} 控制，得到 \eqref{eq:r2:18}。论证不要求截面的零点横截。证毕。

\section{O10 的必要结论：一个普适正量的反自对偶 Weyl 集中}\label{ch2:o10-ux7684ux5fc5ux8981ux7ed3ux8bbaux4e00ux4e2aux666eux9002ux6b63ux91cfux7684ux53cdux81eaux5bf9ux5076-weyl-ux96c6ux4e2d}

定义

\[
 \mathfrak C_-:=\lim_{T\to\infty}\limsup_j
 \frac1{V_j}\int_{\{|W^-_{g_j}|>T\}}|W^-_{g_j}|^2\,dV_j.
 \tag{2.22}\label{eq:r2:22}
\]

幅度与范数都用 O10 的 \texttt{s=-4} 规范。这个量不是某个任选局部气泡的能量；它在原闭流形上计量，自动除以原总体积。

\subsection{定理 5.1}\label{ch2:ux5b9aux7406-5.1}

在 \eqref{eq:r2:1}--\eqref{eq:r2:3} 及原 O10 拓扑输入下，存在普适 \(\epsilon_W>0\)，使任意假想零端序列满足

\[
 \boxed{\mathfrak C_-\ge\epsilon_W>0.}
 \tag{2.23}\label{eq:r2:23}
\]

因此 \(|W^-|^2\) 一致可积的 O10 零端序列不存在。

\subsection{证明}\label{ch2:ux8bc1ux660e-3}

在 \texttt{E\_j=S\^{}-\_j\textbackslash{}otimes\ L\_j} 上，曲率分解写成

\[
 F^{E_j}=F_{0,j}+\mathscr L(W^-_{g_j})+\mathcal E_j,
 \qquad |F_{0,j}|\le C,
 \qquad\|\mathcal E_j\|_{2,\mathrm{av}}\to0.
 \tag{2.24}\label{eq:r2:24}
\]

这里 \texttt{F\_0} 保留标量块、实际有界 \texttt{F\_A\^{}+} 和有界 \texttt{F\_A\^{}-} 模型。\texttt{mathscr\ L} 是四维固定的曲率表示映射。余项包含无迹 Ricci 块及 \texttt{F\_A\^{}-\ -\ iK}，均由原实际预算控制。这只是 Levi-Civita 手性曲率分解；不对模型求导。

若 \texttt{mathfrak\ C\_-} 有限，则对某一固定 T，尾部 limsup 有限，因而 \texttt{sup\_j\ fint\ \textbar{}W\^{}-\textbar{}\^{}2}（去掉有限多个指标后）有限。由 \eqref{eq:r2:24}，引理3.1对命题2.1的 u\_j 可用。若 \texttt{mathfrak\ C\_-} 为无穷，\eqref{eq:r2:23}自动成立。

点态由 \eqref{eq:r2:24} 可证明，对充分大 T，

\[
 \fint_{\{|F^{E_j}|>T\}}|F^{E_j}|^2
 \le C\fint_{\{|W^-|>T/C\}}|W^-|^2
       +C\fint|\mathcal E_j|^2.
\]

因而 \(\mathfrak C_{F^E}\le C\mathfrak C_-\)，其中 C 不依赖 Euler 密度上界。由命题4.1，

\[
 \limsup_j\frac{|c_2(E_j)[M_j]|}{V_j}
 \le C_*\mathfrak C_-.
 \tag{2.25}\label{eq:r2:25}
\]

另一方面，原稿的特征数恒等式为

\[
 c_2(E_j)[M_j]
 =5\chi(M_j)+6\operatorname{sign}(M_j)
 +\frac94(a_j^2-T_{+,j}).
 \tag{2.26}\label{eq:r2:26}
\]

Chern--Gauss--Bonnet 与 signature 给

\[
 \frac{8\pi^2(5\chi+6\operatorname{sign})}{V}
 =9\fint|W^+|^2+\fint|W^-|^2
       +\frac{10}{3}-\frac52\fint|\operatorname{Ric}^\circ|^2.
\]

由于 \texttt{a\_j\^{}2\textgreater{}=T\_\{+,j\}}，

\[
 \liminf_j\frac{c_2(E_j)[M_j]}{V_j}
 \ge d_*:=\frac{5}{12\pi^2}>0.
 \tag{2.27}\label{eq:r2:27}
\]

\eqref{eq:r2:25}--\eqref{eq:r2:27} 给出 \(\mathfrak C_-\ge d_*/C_*\)。例如取

\[
 \epsilon_W:=\frac{d_*}{2C_*}>0
\]

即可得到 \eqref{eq:r2:23}。所有进入尾部系数 \texttt{C\_*} 的常数，只依赖维数、秩、固定幅度及 cutoff，故门槛是普适的。证毕。

\subsection{解释}\label{ch2:ux89e3ux91ca}

最新 R2 使实际扭率接近一个非零常长度，并接近一个粗 Laplacian 特征截面。我们利用它构造同丛的有界 u\_j。\textbf{若相反手性曲率不集中，则 u\_j 的小值区不能承载正密度的 c2。} 而原拓扑与曲率恒等式强迫 c2 有正密度。因此高幅度 W- 集中是逃逸不可缺少的结构。

这里没有重新加入原有有符号恒等式中已经相消的 \texttt{20V/3}，也没有证明旧的平方差本身为正。使用的是同一丛的独立 Chern-class 测试，依赖关系完全展开在上文。

\section{真正扩大 O10 范围：包含临界 3/8，并越过它一个普适正量}\label{ch2:ux771fux6b63ux6269ux5927-o10-ux8303ux56f4ux5305ux542bux4e34ux754c-38ux5e76ux8d8aux8fc7ux5b83ux4e00ux4e2aux666eux9002ux6b63ux91cf}

原稿已经证明：若 \(\chi(M_j)/A_j\le B\)，则沿零端序列

\[
 \limsup_j\fint|W^-|^2\le\frac{16B}{9}-\frac23.
 \tag{2.28}\label{eq:r2:28}
\]

原来仅由非负性只能排除 \texttt{B\textless{}3/8}。现在用定理5.1，

\[
 \epsilon_W\le\mathfrak C_-
 \le\limsup_j\fint|W^-|^2
 \le\frac{16B}{9}-\frac23.
 \tag{2.29}\label{eq:r2:29}
\]

取

\[
 \boxed{\delta_{\rm top}=\frac9{32}\epsilon_W>0.}
 \tag{2.30}\label{eq:r2:30}
\]

若 \(B\le3/8+\delta_{\rm top}\)，\eqref{eq:r2:29}右端不超过 \texttt{epsilon\_W/2}，矛盾。

\subsection{推论 6.1：全参数版本}\label{ch2:ux63a8ux8bba-6.1ux5168ux53c2ux6570ux7248ux672c}

在原实际供给、\(p\ge A>0\)、\(a^2\ge T_+\) 及

\[
 \frac{\chi(M)}A\le\frac38+\delta_{\rm top}
\]

下，存在正的统一 c，使

\[
 \boxed{\fint|\operatorname{Ric}+g|^2+D_\beta/V
 \ge c\frac{(\beta-2)^2}{\beta^2},\qquad\beta\ge2.}
 \tag{2.31}\label{eq:r2:31}
\]

\textbf{证明全参数与原度量范围。} 若不存在 c，可选比值趋零的序列。对于 beta 在任意固定紧区间内且 beta\textgreater2，既有有限参数估计排除该比值趋零；beta=2的目标右端为零。因此 beta 必须趋无穷，且 kappa 趋零。由原稿的精确共形比较，转成 \texttt{s=-4} 的零端序列，类平方、p、A和chi/A不变。定理5.1与\eqref{eq:r2:28}排除它。证毕。

这将此前的严格低密度范围扩展到临界点及一个正邻域。\textbf{不是证明任意B的完整O10。} 在更大的B范围，\eqref{eq:r2:29}允许足够多的相反手性集中，仍需更强的集中限制或真实闭流形构造。

原O11平衡族在充分大指标本来已经满足 \(\chi/A<1/3\)，所以不能把本文新阈值当作首次解决该族端点。新增范围是临界 \texttt{3/8} 及其正邻域；新增机制是实际场到同丛辅助截面的转移和曲率尾部的拓扑检测。

\begin{thebibliography}{P99}
\small
\bibitem[R02]{R02} 本会话研究原稿，\emph{同丛有界截面与第二陈类的曲率集中}，2026-09-29。\path{O10_Resolvent_Chern_Concentration.md}。

\bibitem[P03]{Strong} 所提供网页续研，\emph{O10 Strong Current}，2026-09-29，R1--R2 及同解曲率拟合。

\end{thebibliography}
\end{document}
